\Article{11}
Since the determination of the most plausible errors is independent of the coefficients \(l,l',l''\) etc., it clearly provides the most convenient preparation for whatever use one may wish to make of the observations. Furthermore, it is clear that this task does not require \textit{indefinite} elimination, or knowledge of the coefficients \([\alpha\alpha],[\alpha\beta]\) etc.; all that is required is to derive the auxiliary quantities \(A,B,C\) etc., which we shall hereafter call the \textit{correlates} of the conditional equations \(X=0,\ Y=0,\ Z=0\) etc., from equations (12) by definite elimination and substitute them in formulas (10).

Although this method leaves nothing to be desired whenever only the most plausible values of quantities depending on the observations are required, the situation nevertheless seems to be different whenever the weight of a determination is also desired. For this purpose, depending on which of the four expressions given above one wishes to use, knowledge of the quantities \(L,L',L''\) etc., or at least of the quantities \(x^0,y^0,z^0\) etc., seems necessary. It will therefore be useful to investigate elimination more closely, which will open to us an easier way of finding the weights as well.

\Article{12}
The relation among the quantities occurring in this investigation is considerably clarified by introducing the function of the second degree in the indeterminates
\[
\begin{aligned}&[aa]xx+2[ab]xy+2[ac]xz+\text{etc.}\\&\qquad+[bb]yy+2[bc]yz+\text{etc.}+[cc]zz+\text{etc.}\end{aligned}
\]
which we shall denote by \(T\). First, it is immediately evident that this function becomes
\SourceLeaf{074}
\[
\begin{aligned}&\frac{(ax+by+cz+\text{etc.})^2}{p}+\frac{(a'x+b'y+c'z+\text{etc.})^2}{p'}\\&\qquad+\frac{(a''x+b''y+c''z+\text{etc.})^2}{p''}+\text{etc.}\quad(15)\end{aligned}
\]

Further, it is clear that
\[
T=x\xi+y\eta+z\zeta+\text{etc.}\qquad(16)
\]
and if here \(x,y,z\) etc. are again expressed in terms of \(\xi,\eta,\zeta\) etc. with the aid of equations (7),
\[
\begin{aligned}
T={}&[\alpha\alpha]\xi\xi+2[\alpha\beta]\xi\eta+2[\alpha\gamma]\xi\zeta+\text{etc.}\\&+[\beta\beta]\eta\eta+2[\beta\gamma]\eta\zeta+\text{etc.}\\
&+[\gamma\gamma]\zeta\zeta+\text{etc.}
\end{aligned}
\]

The theory developed above contains two systems of determinate values of the quantities \(x,y,z\) etc. and \(\xi,\eta,\zeta\) etc.; to the first, in which \(x=x^0,y=y^0,z=z^0\) etc. and \(\xi=-[al],\ \eta=-[bl],\ \zeta=-[cl]\) etc., there will correspond the following value of \(T\):
\[
T=[ll]-\frac1P
\]
This is immediately evident either from the third expression for the weight \(P\), compared with equation (16), or from the fourth; to the second, in which \(x=A,y=B,z=C\) etc. and \(\xi=\mathfrak A,\ \eta=\mathfrak B,\ \zeta=\mathfrak C\) etc., there corresponds the value \(T=S\), as is evident either from formulas (10) and (15), or from (14) and (16).

\clearpage
